上一节确定了各对象的保存位置。现在沿一项组合购买说明责任如何穿过这些位置：库存决定先授予付款资格，原 C/T 驱动外部付款，成功结果建立原 EV，消费事务接受初次履约，独立扫描器把发货意图送到效果平台。同步 HTTP 响应观察这条链经过的某个位置，后续推进依靠已经保存的责任。

\subsection{资源决定位于实际发送以前}

组合商品同时需要库存分区 b2 中的 SKU-A 和 b3 中的 SKU-B。购买权威可以先保存 $A_J$ 以及 C/T，把 \code{reservation\_ready} 置为 \code{false}。这样即使调用方断开，服务端也有明确的购买与库存协调责任。C 已经存在，说明责任已登记，真正付款的发送资格来自另一个持久事实。

\CSixteenFigure{ch16-05-gate}{C/T 先承担持久责任，联合预留的 commit 决定再授予付款发送资格}{fig:ch16-gate}

协调记录 TX 保存参与者、子操作身份、原数量与投票。参与者的 \code{prepare} 把数量从可用量转入准备占用，同一子操作重试取得原投票。协调者取得全部 yes 以后，在自己的复制权威中保存不可改的 commit 决定，恢复工作者据该决定完成各参与者。购买权威在核对这一决定与完成状态以后把 \code{ready} 置为 \code{true}。

每一条实际付款路径都经过同一资格检查。扫描器在读取待办时筛选 \code{ready}，权威 \code{claim} 再次在自己的仲裁点核验 \code{ready}。若协调者刚保存 C 便退出，新的支付工作者也只能看到 \code{blocked} 责任。若 commit 决定已保存而通知尚未完成，协调恢复者从原 TX 继续。库存协议的内部原子决定与外部支付按不同保存域接续，第十九章会给出投票、决定和恢复的完整结构。

跨信任边界时，参与者还要取得可验证的决定来源。受信工作者转发决定适用于本章明确的管理角色模型，跨组织接口则可以由参与者查询协调权威或核验签名决定。决定引用原 TX 和参与者集合，付款资格也要对回原 C 所属的原 TX，接收者据此判断来件中的 commit 是否具有授权依据。

\subsection{付款未知如何持续推进}

\code{ready} 成立以后，工作者领取原 T，使用原 \code{scope/pay-C} 向平台发送。平台可以尚未收到、已受理 \code{pending}，或已经完成 X1 但响应在网络中丢失。工作者通过相同身份继续查询或安全重试。原 C 记录最近一次观察、下一次工作时间与重试预算，重启后的扫描器据这些字段枚举责任。

执行者退出后，平台已经接受的操作可以继续完成。旧工作者发送的 HTTP 请求可以继续在网络中到达，新工作者也可能同时用原键重试。平台把这些到达汇入原操作，购买权威只接受符合当前写资格并核验完整关联的证据。旧 \code{owner} 本地写被拒以后，原 C 继续保存，后继工作者从平台找回同一结果。

取得成功 X1 以后，权威的一次提交共同固定 \code{provider\_result}、\code{payment\_result}、\code{O:b1:1}、\code{EV:b1:1} 和未发布责任。任何读到这次成功决定的恢复者都能重建同一个订单与事件。返回给调用方的响应若丢失，查询从这些原关联继续取得结果，后台传播也已经拥有自己的扫描入口。

\subsection{原事件跨过两个提交边界}

下面是 \code{EV:b1:1} 的完整报文。\code{event\_id} 标识业务事件，\code{event\_type} 说明它是购买付款成功，\code{event\_version} 指定载荷解释方式，\code{occurred\_at} 属于原事件内容。重复发布沿用整个报文，恢复时保留原发生时间。

\begin{BookCode}
{
  "event_id": "EV:b1:1",
  "event_type": "purchase.paid",
  "event_version": 1,
  "occurred_at": "2026-10-18T01:00:00Z",
  "payload": {
    "purchase_id": "P:b1:1",
    "charge_id": "C:b1:1",
    "order_id": "O:b1:1",
    "provider_charge_id": "X1",
    "currency": "CNY",
    "amount_minor": 3000
  }
}
\end{BookCode}

时间值在这个假设中固定。生产系统的事件时间来自受理结果所采用的时间来源，在结果提交时保存，排序和去重继续依赖相应身份与日志位置。同一秒可以产生许多事件，它们用各自的 EV 区分。

\CSixteenFigure{ch16-06-boundaries}{结果提交、发布确认、消费提交和独立发货分别拥有持久边界}{fig:ch16-boundaries}

图~\ref{fig:ch16-boundaries} 左侧的一次权威提交把业务结果和发布责任共同保存。发布器报告消息域是否已经承担该 EV 的保存责任。右侧的消费事务再把该事件转成初次履约授权。两个业务提交之间允许消息被重复发送、重复交付，恢复者依据两端已经保存的原关联处理这些重复。

消费事务的处理顺序是：解析已支持版本，查 I 并比较完整原事件，定位 B，核对原订单载荷，取得或创建实际 F，然后共同保存 I、B、F、dispatch 意图与原接收回执。任何冲突都回滚这笔事务。事务提交以后才能确认本次 broker 投递或提交消费位移。

\subsection{同一份回执推进进度并终结交接责任}
\label{sec:ch16-handoff}

发布确认与消费提交之间可能隔着很长的停机期。设消息域已确认 \code{EV:b1:1}，购买权威也已保存 \code{published}，随后消费者停机超过队列 TTL 或日志保留期，消息被清除。此时购买已经 paid，却没有 I、F501 或 dispatch 意图。仅扫描未发布事件会漏过这项购买。消息域通过 TTL 和保留策略限制资源，购买权威则根据尚未取得的交接证据继续承担恢复责任~\cite{rabbitmq-queues,kafka-design}。

为此，Outbox 另保存 \code{handoff\_confirmed} 和下一次核查时间。恢复扫描枚举已有成功结果及原 EV、却尚无交接证据的记录，先到履约域按 EV 查询持久回执。查空时按退避安排重新发布完整原报文，发送过程取得确认后持久更新 \code{published} 及本次确认时间。扫描由已保存责任周期触发，即使队列已经清空、系统此后没有新购买，也会再次找到这一项工作。

回执在 I/B/F/dispatch 意图的同一事务中形成，保存 \code{event\_id}、原订单 O、实际 F501 和接受状态。受控查询同时返回原 I 中的完整 EV、履约来源域及来源记录身份，并核对 B 与 dispatch 意图的已提交关联，接口版本固定这些字段的解释。履约入口只能从已提交 I 返回这份回执与证据。重投同 EV 时先比较整个原报文，再返回原回执；不同 EV 若汇入相同 B，则核对同一订单的授权输入，记录本 EV 与原 F 的关联。履约库中的关联约束把 I、B、F 和意图连接起来，回执不能由尚未提交的临时 F 构造。

交接恢复者通过已授权的履约接口查询，购买权威核验来源域及其身份凭据，再逐项核对接口版本、完整原事件、O、F、业务关联和来源记录身份。仅有一条声称 \code{fulfilled} 的消息，或只比较 EV 字符串，都不足以改变责任。核验通过以后，在一次权威提交中保存原回执及来源证据，登记 \code{handoff\_confirmed=true}，并把查询进度推进为 \code{fulfillment\_accepted}。相同回执重到返回已有事实，异事件内容、异订单或不相容 F 保存为冲突并保留恢复责任。

\CSixteenFigure{ch16-22-handoff}{发布确认先持久保存，核验后的消费回执共同推进履约进度并停止补发}{fig:ch16-handoff}

图~\ref{fig:ch16-handoff} 把购买侧的迁移分成三个状态。结果提交形成 \code{unpublished}，消息确认持久写回后进入 \code{published}，消费回执持久写回后才进入 \code{handoff\_confirmed}。消费者可能很快提交，回执比发布标记更早到达；恢复者先保留可重新查询的回执线索，待发布确认完成持久写回后再提交交接。若原发布确认丢失，就重发同 EV 取得确认，再核读原 I。交接入口对已确认状态作单调更新，迟到的发布标记或重试调度不能把它退回待补发。

消费库已经提交而回执响应丢失时，恢复者重查原 I。购买权威保存回执以前退出时，补发最多产生另一份同 EV，消费端恢复同一 F501。权威保存回执以后退出时，后继扫描器读到已确认状态便停止补发。此时周期 dispatch 扫描已从消费事务接过实际发货责任，履约域的备份、复制与灾难恢复必须共同保存 I/B/F、原回执和未完成意图。交接完成后的购买侧补发已停止，履约库回退也要从自己的恢复来源找回这些责任。

\subsection{dispatch 从已提交意图独立恢复}

独立 dispatch 扫描器读取已提交且尚无发货返回的意图，以 \code{dispatch:initial:O:b1:1} 和原 body 调用发货平台。平台在自己的事务内保存原输入、原返回和效果。若扫描器读到成功返回以后退出，本地意图尚未标为完成，后继扫描器便沿同一键取回原返回。

这条恢复的触发源是消费库中的未完成意图。扫描器启动时先枚举待办，之后按下一次执行时间周期扫描，并可由提交通知缩短等待。消息消费者已经确认 EV、原队列已经清空，甚至系统此后再没有一项新购买，这份意图也会被重新找到。相反，若只在收到下一条消息时顺便检查上一次发货，最后一条消息触发的发送一旦失联，就可能长期没有执行机会。

独立扫描让消息接收事务很快结束，外部发货延迟和退避也能使用自己的并发预算。代价落在额外意图写入、扫描索引、领取仲裁和重复查询上。扫描间隔越长，空闲成本越低，恢复延迟也越长。通知可以降低正常延迟，周期扫描负责找回遗漏通知。图~\ref{fig:ch16-dispatch} 中，业务事务、领取代次和平台原键分别承担持久责任、本地写资格与外部效果合流，运维按意图年龄发现停滞。

\CSixteenFigure{ch16-17-dispatch}{发货恢复由已提交意图独立触发，通知缩短延迟，周期扫描补回遗漏}{fig:ch16-dispatch}

\subsection{确定冲突与瞬时失败的处理}

消息可能由于数据库暂时不可用而处理失败，也可能携带同 EV 异载荷的确定冲突。暂时失败保留当前投递责任并安排重试，确定冲突保存原始字节、来源位置、校验原因和处置状态，暂停受影响的消费范围，交给受控核查入口。

例如 Kafka 分区中 offset 0 是 3000 分的 \code{EV:b1:1}，offset 1 又用同一个 EV 声称 3100 分，offset 2 是另一项正常事件。消费者在 offset 1 发现冲突以后，若跳到 offset 2 成功处理，再提交下一位移 3，就会宣称此前前缀已经处理完，使损坏记录被越过。

本章采用明确的停止策略：消费端持久登记阻断记录，当前处理不提交位移，工作者结束。重启先检查未处置阻断记录，保持停止。较大系统可以按分区暂停，让其他独立分区继续，但提交水位只越过已处置的连续前缀。运维选择修复原数据、验证后隔离或接受一项明确的业务拒绝时，保存处置证据和操作者，再推进该前缀。

瞬时错误则回到原位置。Kafka 的 \code{poll} 已经推进本地读取位置，处理失败以后要 \code{seek} 回原记录，或关闭消费者后从已提交位移恢复。RabbitMQ 的未确认投递可以保留，也可以以明确重入队策略重试。重试采用退避和有界并发，避免一条暂时失败的记录占满所有恢复资源。

死信队列保存暂时退出正常消费路径的消息，供隔离、诊断和受控重放。常见原因包括消费者明确拒绝且要求不重入原队列、消息过期、队列长度限制或达到交付次数上限。RabbitMQ 的死信交换机接收这类转交，再按绑定把消息路由到专门的队列。因此，配置一项死信路由时，要同时指定源队列的转交政策、死信交换机和目标队列的保存方式~\cite{rabbitmq-quorum-queues,rabbitmq-queues}。

图~\ref{fig:ch16-dead-letter} 中，一条确定冲突消息被隔离以后，原始字节、原位置和处置原因继续可查。将消息送入死信域需要先确认该域已经持久接受，再按明确策略处置原投递。RabbitMQ quorum queue 的至少一次死信转移显式设置 \code{dead-letter-strategy=at-least-once}、\code{overflow=reject-publish} 和死信交换机，并启用所需 feature flag。内部消费者获得目标确认以前保留源消息，目标不可用会形成积压和反压。默认策略采用至多一次转交，存在转交丢失窗口~\cite{rabbitmq-quorum-queues}\footnote{死信转交也可能重复，消费者按原 EV 和处置身份接收。积压的死信占用源端或目标端资源，其额度和告警一同纳入恢复预算。}。

消息成功进入死信队列，只确定隔离副本已有保存位置。业务冲突的原证据、该订单是否已经完成履约，以及允许推进的连续位置，共同决定后续处置。核查者若确认 3100 分报文损坏，可保存业务拒绝及隔离凭据，让原 3000 分的 I/B/F 保持原状。需要重放时沿原 EV 重放经过核准的原输入，使已有成功通过原关系恢复。

\CSixteenFigure{ch16-18-dead-letter}{瞬时错误回原位置重试，确定冲突经持久隔离和业务处置后才推进前缀}{fig:ch16-dead-letter}

\subsection{退款作为独立的跨域责任}

原 \code{C:b1:1} 成功收款 3000 分以后，每项退款取得自己的 $S_{\rm RF}$ 和 RF。来源 X1 确认可退的原收款，RF 识别这一项退款。两项各退 1000 分可以合法并存，它们各自的网络重试只恢复自己的原退款结果。退款可以由另一组工作者处理，额度仲裁放在原收款所属权威中。沿用第十五章的符号，设原收款额为 $A$，已经成功退款为 $R_{\rm succ}$，尚未排除成功的退款占用为 $R_{\rm pend}$，则可接受的新退款金额 $x$ 满足
\[
R_{\rm succ}+R_{\rm pend}+x\leq A,\qquad A=3000\text{。}
\]
未知退款先占用 $R_{\rm pend}$，以免后续退款抢走同一份额度。

\CSixteenFigure{ch16-07-refunds}{退款先在原收款权威占额，成功把未知占用转成成功金额}{fig:ch16-refunds}

权威先按新 RF 仲裁原输入。已有 RF 的恢复取得原关联，新 RF 才按当前成功与未知占用检查额度，并共同保存 RF、原平台退款身份和占额。成功使 $R_{\rm pend}$ 减少 $x$、$R_{\rm succ}$ 增加 $x$，总占用保持。足够的最终拒绝才释放 $R_{\rm pend}$，普通超时与查空保留原占用。

图~\ref{fig:ch16-refunds} 中的 \code{RF:b1:4} 先占 2000 分，随后一项 1200 分退款超过剩余 1000 分，当前首次受理被拒绝。平台已完成 2000 分而本地保存以前工作者退出，原 RF 的恢复把相应金额从 $R_{\rm pend}$ 转入 $R_{\rm succ}$。新的 RF 再退 1000 分，合计达到 3000。退款工作者、付款工作者与发布器都从各自持久责任扫描，领取代次变化以后继续使用原外部身份。
